\documentclass[11pt,a4paper]{article}

\usepackage[utf8]{inputenc}
\usepackage[T1]{fontenc}
\usepackage{lmodern}
\usepackage[english]{babel}
\usepackage{geometry}
\geometry{a4paper, left=2.5cm, right=2.5cm, top=2.5cm, bottom=2.5cm}
\usepackage{graphicx}
\usepackage{microtype}
\usepackage{amsmath,amssymb,amsfonts}
\usepackage{booktabs}
\usepackage{multirow}
\usepackage{array}
\usepackage{caption}
\usepackage{subcaption}
\usepackage{xcolor}
\usepackage{hyperref}
\usepackage{siunitx}

\graphicspath{{figures/}}
\hypersetup{colorlinks=true, linkcolor=blue, citecolor=blue, urlcolor=blue}
\sisetup{detect-all}

\newcommand{\rps}{\ensuremath{\mathrm{RPS}}}
\newcommand{\revps}{\ensuremath{\mathrm{rev}/\mathrm{s}}}
\newcommand{\rsq}{\ensuremath{R^{2}}}

\title{\textbf{DREGON-LM-V2: A Harder Evaluation for Multi-Rotor \rps{} Prediction}\\
  \large Cross-evaluation of old checkpoints and new V3 models}
\author{Dmitrii Mukhutdinov}
\date{\today}

\begin{document}
\maketitle

% ============================================================================
\section{Why a new dataset?}
% ============================================================================

The original DREGON-LM dataset (Papers~1 and~2) is not great. Three reasons:

\begin{enumerate}
  \item \textbf{Train/validation overlap.} Both splits come from the same
    recordings, so a model can just memorise room acoustics and motor
    signatures instead of actually learning to predict RPS.
  \item \textbf{Short clips.} 0.82~seconds $\approx$ 26 STFT frames. Not
    enough temporal context for recurrent heads to do anything useful.
  \item \textbf{One microphone only.} Eight microphones exist in DREGON; we
    used one. Spatial diversity was ignored.
\end{enumerate}

DREGON-LM-V2 fixes all of this. We re-train on V2 (call them \emph{V3}
checkpoints) and cross-evaluate every model on both the old and the new
validation set. If a model was only good because it memorised the old split,
we will know.

% ============================================================================
\section{What changed in V2}
% ============================================================================

\paragraph{Sources.} Same ingredients: DREGON (eight channels, four motors)
and Librispeech. The recipe is different:
\begin{itemize}
  \item Pick one of the eight microphones per sample. Spatial diversity.
  \item Mix speech at $-30$ to $0$~dB SNR.
  \item Use \emph{commanded} motor speeds (flight-controller setpoints)
    instead of measured telemetry. Cleaner ground truth, no alignment headaches.
\end{itemize}

\paragraph{Duration.} 3 seconds (48000 samples at 16~kHz) = $\sim$94 STFT
frames. More than 3$\times$ the old length. Recurrent heads can actually
breathe now.

\paragraph{Train/valid split.} Recording-level, zero overlap. Train on 6
\texttt{in\_flight\_nosource} recordings; validate on two completely
unseen \texttt{free-flight} recordings. No memorisation possible.

\paragraph{Synthetic clips.} About 20\% of training samples are synthesised
by mixing individual motor recordings at constant speed. The idea was sound:
force the model to learn actual pitch tracking instead of just predicting
rough in-flight statistics. The execution was heavy-handed---20\% is too
much, and the model learns to default to ``predict the mean'' whenever it
is uncertain. More on this later.

\paragraph{Stats.} 6000 train + 600 valid, 2.2~GB. Old dataset stays where
it is; we are not touching it.

% ============================================================================
\section{Training setup}
% ============================================================================

Two new models trained on V2 (the \textbf{V3} checkpoints):

\begin{itemize}
  \item \textbf{V3 SimpleConv} (baseline): 5 encoder blocks, global average
    pool + 1-D conv head. 538~K params.
  \item \textbf{V3 BiGRU-v2} (best from the architecture sweep): 6 blocks,
    squeeze--excitation, BiGRU head. 1.44~M params.
\end{itemize}

Standard recipe: STFT $(N=2048, H=512)$ at 16~kHz, AdamW ($10^{-3}$, weight
decay $10^{-4}$), ReduceLROnPlateau, early stopping (patience~30).
\textbf{PIT loss} because the four rotor outputs are unordered---we try all
$4! = 24$ permutations and backprop the cheapest one. Batch size 96 (up from
16 in the first V2 run) because why not.

We also keep the \textbf{OLD} checkpoints (trained on the original
DREGON-LM) around for comparison: OLD SimpleConv and OLD BiGRU-v2.

% ============================================================================
\section{Cross-evaluation results}
% ============================================================================

Every model is tested on \emph{both} the old validation set (1-second clips,
measured RPS, train/val overlap) and the V2 validation set (3-second clips,
command RPS, zero overlap). Metrics are \textbf{PIT-MSE}.

Table~\ref{tab:cross} and Figure~\ref{fig:cross} tell the story.

\begin{table}[ht]
  \centering
  \caption{Cross-evaluation PIT-MSE [(rev/s)$^2$] and degradation.}
  \label{tab:cross}
  \small
  \begin{tabular}{lcccc}
\toprule
Model & OLD valid & V2 valid & Degradation & V2 MAE \\
\midrule
OLD SimpleConv & 5.2 & 331.9 & 63$\times$ & 7.68 \\
OLD BiGRU-v2 & 2.7 & 327.3 & 123$\times$ & 6.29 \\
V3 SimpleConv & 66.8 & 148.1 & 2.2$\times$ & 10.57 \\
V3 BiGRU-v2 & 15.3 & 71.1 & 4.7$\times$ & 4.34 \\
\bottomrule
\end{tabular}

\end{table}

\begin{figure}[ht]
  \centering
  \includegraphics[width=\textwidth]{fig_cross_eval}
  \caption{Old models look amazing on the old set (blue) and die on V2
    (orange). V3 models are worse on the old set but far more robust on V2.}
  \label{fig:cross}
\end{figure}

\textbf{Old models were cheating.} OLD BiGRU-v2 scores PIT-MSE 2.7 on the
old validation set---then explodes to 327 on V2. That is a
$\mathbf{123\times}$ degradation. OLD SimpleConv goes from 5.2 to 331.9
($63\times$). The old validation split is not a generalisation test; it is
a memorisation test. The models learned recording-specific signatures, not
RPS prediction.

\textbf{V3 models actually generalise.} V3 BiGRU-v2 degrades only
$4.7\times$ (15.3 $\to$ 71.1) and V3 SimpleConv only $2.2\times$
(66.8 $\to$ 148.1). The absolute V2 numbers are still higher than the old
inflated ones, but the \emph{relative} robustness is what matters.

\begin{figure}[ht]
  \centering
  \includegraphics[width=0.7\textwidth]{fig_degradation}
  \caption{Degradation factor: old models collapse; V3 models degrade
    modestly.}
  \label{fig:degradation}
\end{figure}

% ============================================================================
\section{Why old models fail on V2}
% ============================================================================

Three reasons, all tied to the old dataset being too easy:

\textbf{Memorisation.} Old train and validation come from the same
recordings. A model can just learn the room, the microphone position, and
the characteristic motor speeds of each recording. With 0.82-second clips,
predicting the per-recording mean RPS is already pretty good.

V2 removes all three cheats: different recordings for validation, all eight
microphones, and 3-second clips with actual temporal trajectories.

\textbf{Command vs measured RPS.} V2 uses flight-controller setpoints
instead of measured telemetry. Setpoints are smoother but noisier around
transitions, and they lag the physical rotor speed by a variable amount.
Harder target.

\textbf{Temporal context.} 3~seconds = $\sim$94 STFT frames. The model
must follow a trajectory, not guess a scalar. OLD BiGRU-v2 (trained on
1-second clips) scores 327 on V2; V3 BiGRU-v2 (trained on 3-second clips)
scores 71. Context matters.

% ============================================================================
\section{PIT vs standard MSE: does the model know which rotor is which?}
% ============================================================================

PIT loss hides the permutation problem during training. But we can expose it
by comparing PIT-MSE (best permutation) against standard MSE (fixed channel
order).

\begin{figure}[ht]
  \centering
  \includegraphics[width=0.65\textwidth]{fig_pit_std_gap}
  \caption{Std-MSE excess over PIT-MSE on V2 valid. Big gap = the model
    has no idea which rotor is which.}
  \label{fig:gap}
\end{figure}

V3 SimpleConv: 38\% gap (PIT 148 vs Std 205). Without a recurrent head,
the CNN cannot keep rotors in consistent channels. V3 BiGRU-v2: only 4\%
gap (71 vs 74). The BiGRU learns a stable temporal ordering. This confirms
the architecture-sweep finding: temporal heads are not optional for
multi-rotor outputs.

% ============================================================================
\section{Per-channel analysis on V2}
% ============================================================================

Old dataset always uses channel~0. Do old models get a ``home-field''
advantage on ch0?

\begin{table}[ht]
  \centering
  \caption{PIT-MSE on V2 valid, split by microphone channel.}
  \label{tab:channel}
  \small
  \begin{tabular}{lccc}
\toprule
Model & ch0 & ch1--7 & ch1--7/ch0 \\
\midrule
OLD SimpleConv & 293.6 & 338.0 & 1.15$\times$ \\
OLD BiGRU-v2 & 247.2 & 340.1 & 1.38$\times$ \\
V3 SimpleConv & 126.9 & 151.5 & 1.19$\times$ \\
V3 BiGRU-v2 & 66.2 & 71.9 & 1.09$\times$ \\
\bottomrule
\end{tabular}

\end{table}

Yes. OLD models are 15--38\% better on ch0 than on the other seven
channels. V3 models shrink this to 9--19\%, because they trained on all
eight. Even on ch0 alone, V3 BiGRU-v2 beats OLD BiGRU-v2 by $3.7\times$
(66.2 vs 247.2).

% ============================================================================
\section{In-flight recordings with source signals}
% ============================================================================

We also evaluate all four checkpoints on two real in-flight recordings that
contain an additional source signal (speech or white noise) mixed with the
drone noise.  In the raw DREGON files the motor-telemetry stream starts
several seconds after the audio stream and ends before it, so we crop the
audio to the exact motor-timestamp range and interpolate the commanded RPS
to the STFT frame rate of the cropped audio.  Table~\ref{tab:inflight}
reports mean PIT-MSE over the full aligned sequence.

\begin{table}[ht]
  \centering
  \caption{In-flight mean PIT-MSE [(rev/s)$^2$] on aligned full-sequence data.}
  \label{tab:inflight}
  \small
  \begin{tabular}{lcccc}
    \toprule
    & OLD SimpleConv & OLD BiGRU-v2 & V3 SimpleConv & V3 BiGRU-v2 \\
    \midrule
    speech-high  & 224.6 & 199.9 & 274.1 & 277.3 \\
    whitenoise-high & 87.7 & 104.0 & 123.8 & 323.0 \\
    \bottomrule
  \end{tabular}
\end{table}

The old checkpoints---trained on V1, which included in-flight segments and
used measured motor speeds---generalise surprisingly well to in-flight
recordings with source signals (MSE 88--225).  In contrast, the V3
checkpoints---trained exclusively on the static V2 dataset with commanded
speeds---show notably higher error, especially V3 BiGRU-v2 on whitenoise-high
(MSE 323).

The constant-speed synthetic clips were added on purpose: the idea was to
force the model to learn actual pitch tracking instead of just predicting
rough in-flight statistics. The intuition was sound---but 20\% is too
much. The model learns the shortcut: when in doubt, predict the mean.

On a synthetic constant-speed validation sample, V3 BiGRU-v2 prediction
std is $\sim$0.05~rev/s while ground-truth std is $\sim$1.2~rev/s. That
is a 20$\times$ variance suppression. You can see it in
Figures~\ref{fig:inflight-speech-v3} and~\ref{fig:inflight-whitenoise-v3}:
the V3 traces are suspiciously smooth and miss transients.

Fix for V2: keep the synthetic idea, but (1) drop the proportion to maybe
5\%, and (2) time-stretch the individual motor recordings before mixing so
that rotor speeds actually vary within each synthetic clip. Otherwise the
V3 checkpoints look good on paper but fail in the real world.

Full-sequence 3-panel plots below. Top = spectrogram; middle = per-rotor
RPS (dotted = ground truth, solid = prediction); bottom = per-frame MSE.
Error is mostly at spin-up ($t < 5$~s); once rotors stabilise the tracking
is decent.

\begin{figure}[p]
  \centering
  \includegraphics[width=\textwidth]{fig_fullseq_free-flight_speech-high_room1_old_bigru-v2}
  \caption{Full-sequence evaluation on \texttt{free-flight\_speech-high\_room1}
    ($\sim$47~s).  OLD BiGRU-v2 (trained on V1 with in-flight data) tracks the
    four rotors accurately; error is mainly concentrated at spin-up.}
  \label{fig:inflight-speech-old}
\end{figure}

\begin{figure}[p]
  \centering
  \includegraphics[width=\textwidth]{fig_fullseq_free-flight_speech-high_room1_v3_bigru-v2}
  \caption{Same recording, V3 BiGRU-v2 (trained on static V2 only).
    Steady-flight tracking is comparable, but spin-up error is larger
    (mean MSE 277 vs 200).}
  \label{fig:inflight-speech-v3}
\end{figure}

\begin{figure}[p]
  \centering
  \includegraphics[width=\textwidth]{fig_fullseq_free-flight_whitenoise-high_room1_old_bigru-v2}
  \caption{Full-sequence evaluation on \texttt{free-flight\_whitenoise-high\_room1}
    ($\sim$52~s).  OLD BiGRU-v2 again shows strong tracking (mean MSE 104).}
  \label{fig:inflight-whitenoise-old}
\end{figure}

\begin{figure}[p]
  \centering
  \includegraphics[width=\textwidth]{fig_fullseq_free-flight_whitenoise-high_room1_v3_bigru-v2}
  \caption{Same recording, V3 BiGRU-v2.  The model under-estimates rotor
    speeds during spin-up and shows larger drift in steady flight
    (mean MSE 323 vs 104).}
  \label{fig:inflight-whitenoise-v3}
\end{figure}

% ============================================================================
\section{V3 training dynamics}
% ============================================================================

Figure~\ref{fig:v3train} shows the training curves. Both models overfit
after early epochs; early stopping keeps the best validation checkpoint.

\begin{figure}[ht]
  \centering
  \includegraphics[width=\textwidth]{fig_v3_training_curves}
  \caption{Train MSE keeps dropping; validation MSE plateaus. The task is
    genuinely hard, not a convergence issue.}
  \label{fig:v3train}
\end{figure}

V3 SimpleConv: 52 epochs, stops at val PIT-MSE 148.1. V3 BiGRU-v2: 44
epochs, stops at 71.1. Same pattern as the first V2 run: train goes down,
validation stays flat. The task is hard.

% ============================================================================
\section{Why $R^2$ is useless here}
% ============================================================================

$R^2$ comes out as $-10^4$. Not because predictions are bad---because many
V2 samples have near-constant RPS ($\sigma < 1$~rev/s over 3~seconds), so
the baseline variance is basically zero and the ratio explodes. We just
report MSE and MAE.

% ============================================================================
\section{Bottom line}
% ============================================================================

DREGON-LM-V2 is harder and cleaner than V1. The old dataset's validation
scores (MSE 2--5) were a mirage; old checkpoints degrade $63$--$123\times$
on V2. V3 checkpoints degrade only $2.2$--$4.7\times$.

But V2 has a self-inflicted wound. The constant-speed synthetic clips were
added with good intentions---force the model to track pitch, not just
statistics---but at 20\% they become a shortcut. V3 models learn to predict
the mean; on real in-flight data they underperform the old models (MSE
124--323 vs 88--225).

BiGRU-v2 is still the best architecture (PIT-MSE 71.1 on V2 valid vs 148.1
for SimpleConv). The fix is simple: rebuild V2 with fewer constant-speed
synthetic clips, and time-stretch the individual motors so that speeds vary
within each synthetic sample. Then we can trust the V3 numbers.

\end{document}
